
\documentclass{article}
\usepackage{colortbl}
\usepackage{makecell}
\usepackage{multirow}
\usepackage{supertabular}

\begin{document}

\newcounter{utterance}

\centering \large Interaction Transcript for game `wordle\_withcritic', experiment `high\_frequency\_words\_clue\_with\_critic', episode 8 with reflexion{-}llama{-}3.1{-}8b{-}instruct.
\vspace{24pt}

{ \footnotesize  \setcounter{utterance}{1}
\setlength{\tabcolsep}{0pt}
\begin{supertabular}{c@{$\;$}|p{.15\linewidth}@{}p{.15\linewidth}p{.15\linewidth}p{.15\linewidth}p{.15\linewidth}p{.15\linewidth}}
    \# & $\;$A & \multicolumn{4}{c}{Game Master} & $\;\:$B\\
    \hline

    \theutterance \stepcounter{utterance}  
    & & \multicolumn{4}{p{0.6\linewidth}}{
        \cellcolor[rgb]{0.9,0.9,0.9}{
            \makecell[{{p{\linewidth}}}]{
                \texttt{\tiny{[P1$\langle$GM]}}
                \texttt{You have attempted this same task before and did not succeed. These are the plans you wrote after those attempts.} \\
\texttt{Trial \#1: New plan: } \\
\\ 
\texttt{In my last attempt, I made two mistakes. Firstly, I provided multiple guesses in a single message, which is not allowed. Secondly, I provided a guess that was not a valid word. To correct this, I will make sure to provide a single guess in each message, and I will choose a valid word that fits the given clue. In this case, I should have guessed a different TV show that started in 1952, and I should have done it in a single message. For example, I could have guessed "I Love Lucy" in the first message, or I could have guessed a different show like "The Honeymooners" or "The Tonight Show".} \\
\\ 
\texttt{Trial \#2: New plan: } \\
\\ 
\texttt{In my last attempt, I made two mistakes. Firstly, I provided a guess that was not a valid word. Secondly, I did not use the clue effectively to arrive at a new guess. To correct this, I will make sure to choose a valid word that fits the given clue. In this case, I should have used the clue "tv show since 1/14/52" to guess a TV show that started in 1952. I should have also considered the fact that the clue mentions a specific date, which suggests that the TV show is a classic one. Next time, I will use the clue to arrive at a new guess, and I will make sure to choose a valid word that fits the given clue. For example, I could have guessed "The Honeymooners" or "The Tonight Show" in the first message.} \\
\\ 
\texttt{Trial \#3: New plan: } \\
\\ 
\texttt{In my last attempt, I made two mistakes. Firstly, I provided a guess that was not a valid word. Secondly, I did not use the clue effectively to arrive at a new guess. To correct this, I will make sure to choose a valid word that fits the given clue. In this case, I should have used the clue "tv show since 1/14/52" to guess a TV show that started in 1952. I should have also considered the fact that the clue mentions a specific date, which suggests that the TV show is a classic one. Next time, I will use the clue to arrive at a new guess, and I will make sure to choose a valid word that fits the given clue. } \\
\\ 
\texttt{To achieve this, I will follow these steps:} \\
\\ 
\texttt{1. Read the clue carefully and understand its meaning.} \\
\texttt{2. Choose a valid word that fits the given clue.} \\
\texttt{3. Make sure the word is a TV show that started in 1952.} \\
\texttt{4. Provide a single guess in each message.} \\
\texttt{5. Use the clue effectively to arrive at a new guess.} \\
\\ 
\texttt{For example, I could have guessed "The Honeymooners" or "The Tonight Show" in the first message.} \\
\\ 
\\ 
\texttt{You are a language wizard who likes to guess words by using the given rules.} \\
\\ 
\texttt{Welcome to Wordle! You have six attempts to guess the target word, a valid English word of five lowercase letters (a{-}z). Please use the tags "explanation:" and "guess:" to provide a concise explanation for each guess.} \\
\\ 
\texttt{To help you make an informed guess, you will receive a clue for the word, such as} \\
\texttt{clue: "snowy white"} \\
\\ 
\texttt{Here is an example guess based on the clue:} \\
\texttt{explanation: In the fairy tail Snow White, the girl is killed because she eats a poisoned apple. And the word apple has 5 letters.} \\
\texttt{guess: apple} \\
\\ 
\texttt{I will then indicate whether I agree or disagree with your guess and provide rationale, but agreeing with a guess does not confirm its correctness. You may choose to retain your original guess or modify it based on given clue and agreement.} \\
\\ 
\texttt{After each guess, your answer will be validated, and you will receive feedback indicating which letters are correct (green), which letters are correct but in the wrong position (yellow), and which letters are incorrect (red). This feedback can be useful in determining which letters to include or exclude in your next guess.} \\
\\ 
\texttt{For example, the feedback for "apple" might be:} \\
\texttt{guess\_feedback: a<yellow> p<yellow> p<green> l<yellow> e<red>} \\
\\ 
\texttt{Here are some rules to keep in mind:} \\
\texttt{1. Please do not generate any guess feedback or clue on your own} \\
\texttt{2. Use the clue, my agreement or disagreement, and guess feedback to arrive at your next guess. This should be reflected in your explanation as well.} \\
\\ 
\texttt{Do not generate any other text.} \\
\\ 
\texttt{Let's begin with your first guess.} \\
\\ 
\texttt{clue: tv show since 1/14/52} \\
            }
        }
    }
    & \\ \\

    \theutterance \stepcounter{utterance}  
    & \multicolumn{4}{p{0.6\linewidth}}{
        \cellcolor[rgb]{0.9,0.9,0.9}{
            \makecell[{{p{\linewidth}}}]{
                \texttt{\tiny{[P1$\rangle$GM]}}
                \texttt{explanation: The clue mentions a specific date, which suggests that the TV show is a classic one. I will use this information to arrive at a new guess. I will choose a valid word that fits the given clue, which is a TV show that started in 1952. } \\
\texttt{guess: thehoney} \\
            }
        }
    }
    & & \\ \\

    \theutterance \stepcounter{utterance}  
    & & & \multicolumn{2}{p{0.3\linewidth}}{
        \cellcolor[rgb]{0.9,0.9,0.9}{
            \makecell[{{p{\linewidth}}}]{
                \texttt{\tiny{[GM$|$GM]}}
                \texttt{Error: The length of the guessed word is not 5.} \\
            }
        }
    }
    & & \\ \\

    \theutterance \stepcounter{utterance}  
    & & & \multicolumn{2}{p{0.3\linewidth}}{
        \cellcolor[rgb]{0.9,0.9,0.9}{
            \makecell[{{p{\linewidth}}}]{
                \texttt{\tiny{[GM$|$GM]}}
                \texttt{game\_result = ABORT} \\
            }
        }
    }
    & & \\ \\

    \theutterance \stepcounter{utterance}  
    & \multicolumn{4}{p{0.6\linewidth}}{
        \cellcolor[rgb]{0.95,0.95,0.95}{
            \makecell[{{p{\linewidth}}}]{
                \texttt{\tiny{[P2$\rightarrow$P1]}}
                \texttt{Post{-}episode Reflexion (not sent to GM):} \\
\texttt{New plan: } \\
\\ 
\texttt{In my last attempt, I made two mistakes. Firstly, I provided a guess that was not a valid word. Secondly, I did not use the clue effectively to arrive at a new guess. To correct this, I will make sure to choose a valid word that fits the given clue. In this case, I should have used the clue "tv show since 1/14/52" to guess a TV show that started in 1952. I should have also considered the fact that the clue mentions a specific date, which suggests that the TV show is a classic one. Next time, I will use the clue to arrive at a new guess, and I will make sure to choose a valid word that fits the given clue. } \\
\\ 
\texttt{To achieve this, I will follow these steps:} \\
\\ 
\texttt{1. Read the clue carefully and understand its meaning.} \\
\texttt{2. Choose a valid word that fits the given clue.} \\
\texttt{3. Make sure the word is a TV show that started in 1952.} \\
\texttt{4. Provide a single guess in each message.} \\
\texttt{5. Use the clue effectively to arrive at a new guess.} \\
\\ 
\texttt{For example, I could have guessed "The Honeymooners" or "The Tonight Show" in the first message.} \\
            }
        }
    }
    & & \\ \\

\end{supertabular}
}

\end{document}
